% Source: https://artofproblemsolving.com/community/c3h1694352_permutation_and_combination
% Stable identity: aops-topic-1694352-q01
Since $2010=2\cdot3\cdot5\cdot67$ is a product of four distinct primes, each of $a,b,c,d$ is a product of some of these primes, each prime used exactly once. The four numbers are distinct, so at most one of them is $1$.

\textbf{No factor equal to $1$.} Then each of the four numbers is a single prime: $\{2,3,5,67\}$, $1$ way.

\textbf{Exactly one factor is $1$.} Then the four primes are split into three non-empty groups, one group of two primes and two single primes. The pair can be chosen in $\dbinom42=6$ ways, and the resulting numbers ($1$, the product of the pair, and the two other primes) are always distinct: the product of two primes differs from each single prime and from $1$. So $6$ ways.

(Two or more factors equal to $1$ is impossible since the factors are distinct.) Each choice of a set $\{a,b,c,d\}$ gives exactly one ordering $a<b<c<d$, so the number of ways is $1+6=\boxed{7}$.
